\documentclass[12pt]{article}

\input{../../../_assets/preambulo}

\begin{document}

    % 1. Foto de fondo
    % 2. Título
    % 3. Encabezado Izquierdo
    % 4. Color de fondo
    % 5. Coord x del titulo
    % 6. Coord y del titulo
    % 7. Fecha

    
	\input{../../../_assets/portada}
	\portadaExamenFotoDif{../../sapienza.jpg}{Computer Vision\\Examen III}{Computer Vision. Examen III}{MidnightBlue}{-8}{28}{2025-2026}{Joaquín Avilés de la Fuente}

    \begin{description}
        \item[Asignatura] Computer Vision.
        \item[Curso Académico] 2025-2026.
        \item[Grado] Grado en Informática.
        \item[Grupo] Grupo A.
        \item[Profesor] Irene Amerini.
        \item[Descripción] Examen final escrito de Computer Vision.
        \item[Fecha] abril de 2026.
        \item[Duración] 60 min.
    
    \end{description}
    \newpage

    % ------------------------------------

    \begin{observacion}
        The total sum of all the exercises is 32 points, although the maximum mark of the exam is 30.
    \end{observacion}
    
   \begin{ejercicio}[8 points]
    Describe the Sobel filter for edge detection. How does it differ from the Canny edge detector?
    \end{ejercicio}

    \begin{ejercicio}[8 points]
        Explain feature matching and how the RANSAC (Random Sample Consensus) algorithm can be used in conjunction with feature matching to improve the accuracy of homography estimation between two images.
    \end{ejercicio}

    \begin{ejercicio}[8 points]
        \begin{enumerate}
            \item[a.] Explain the concept of fundamental matrix. How is it computed, and what information does it convey about the relationship between two images?
            \item[b.] What is the geometric meaning of the epipoles in the two images? How are they related to the fundamental matrix (algebraically)?
        \end{enumerate}
    \end{ejercicio}

    \begin{ejercicio}[8 points]
        Propose a method to detect keyframes in a video sequence using traditional (non-deep learning) techniques. Explain the features you would extract, how you would measure frame importance or uniqueness, and the criteria for selecting keyframes. Justify each step in your approach. A keyframe in this case could be: a frame that is visually different from surrounding frames (e.g., scene changes, new objects) or a frame where meaningful action begins (e.g., explosion, expression change, speaker switch).
    \end{ejercicio}
\end{document}